\section{Dimensional reconstruction and exterior conservation}
\label{nuclear:11}
\subsection{Two memories recover the centred register}

\subsubsection{Register-reading operators}

In the ordered chart of the register \(K\) from \cref{sec:k-registro,img09:imagen}, let \(V=\RR^{12}\) be its subsequent linear realization and let \(S\) be the permutation \((Sx)_i=x_{i+1}\), with indices modulo \(12\). The readers \(I-S^3\) and \(I-S^4\) are those of the transverse extractor in \cref{img09:imagen}. The nonadic compression in \cref{nuclear:10} provides, for \(j=3,4\),

\begin{equation}
T_j=\frac{8I+S^j}{9},\qquad D_j=\frac{\sqrt8}{9}(S^j-I),
\qquad T_j^*T_j+D_j^*D_j=I.
\label{nuc11:eq:1}
\end{equation}

The partition \(8+1\) and normalization \(9\) come from the nine positions and their collective contrast. Consider the following explicit composition of readers, applied to the generated register:

\begin{equation}
m_0=D_3K,\qquad m_1=D_4T_3K,
\qquad M K=(m_0,m_1).
\label{nuc11:eq:2}
\end{equation}

The second memory acts on the reading produced by the first stage. The complements are archived separately. The complete evolution with re-entry corresponds to the reversible coupling in \cref{nuclear:10}.

\subsubsection{Exact inverse}

Since \((S^3)^4=I\),

\begin{equation}
T_3^{-1}=\frac{512I-64S^3+8S^6-S^9}{455}.
\label{nuc11:eq:3}
\end{equation}

Indeed, \((8I+S^3)(512I-64S^3+8S^6-S^9)=4095I\) and \(4095=9\cdot455\). Hence the memories provide

\begin{equation}
b=-\frac9{\sqrt8}m_0=(I-S^3)K,
\quad c=-\frac9{\sqrt8}T_3^{-1}m_1=(I-S^4)K.
\label{nuc11:eq:4}
\end{equation}

Commutation of \(T_3\) with \(S^4\) justifies the second equality. Define

\[
w_i=c_i-b_{i+1},\quad B_0=0,\quad
B_i=\sum_{j<i}w_j.
\]

The identity \(w_i=K_i-K_{i+1}\) proves that \(B_i=K_0-K_i\). If \(q=\sum_{i=0}^{11}K_i\), we recover exactly

\begin{equation}
K_i=\frac{q+\sum_{j=0}^{11}B_j}{12}-B_i.
\label{nuc11:eq:5}
\end{equation}

This is the inverse \eqref{img09:inversa}, now composed with compression and memory. With \(q=0\), the formula returns \(P_{11}K\), where

\[
P_{11}=I-\frac1{12}\mathbf1\mathbf1^*.
\]

The two memories recover the eleven zero-mean components. The uniform charge is recovered through \(q\). The kernel of \(M\) is precisely the uniform line: vanishing of both differences implies fixation by \(S^3\) and \(S^4\); since \(S=S^4(S^3)^{-1}\), it implies fixation by \(S\).

\subsubsection{Optimal stability bound}

The Gram matrix of the procedure is

\begin{equation}
G=M^*M=D_3^*D_3+T_3^*D_4^*D_4T_3
=I-(T_4T_3)^*(T_4T_3).
\label{nuc11:eq:6}
\end{equation}

In mode \(k\) of \(S\), the eigenvalue of \(T_j^*T_j\) is
\[
 \frac{65+16\cos(2\pi jk/12)}{81}.
\]
Substituting it into the preceding Gram matrix gives, with multiplicities,

\begin{equation}
\operatorname{Spec}G=\left\{
0^{[1]},\ (16/81)^{[2]},\ (8/27)^{[2]},\ (32/81)^{[1]},
(952/2187)^{[4]},\ (1256/2187)^{[2]}\right\}.
\label{nuc11:eq:7}
\end{equation}

The zero mode is the mean. Let \(L\) be the least-squares inverse of \(M\) on \(\mathbf1^\perp\), extended to the register space by

\[
L=(G|_{\mathbf1^\perp})^{-1}M^*.
\]

Then

\begin{equation}
LM=P_{11},\qquad \|L\|=\frac94.
\label{nuc11:eq:8}
\end{equation}

The bound is optimal because modes \(k=3,9\) attain \(16/81\). For perturbed data and reconstruction with \(L\),

\begin{equation}
\|\delta K\|^2\le \frac{|\delta q|^2}{12}
+\frac{81}{16}\bigl(\|\delta m_0\|^2+\|\delta m_1\|^2\bigr).
\label{nuc11:eq:9}
\end{equation}

The uniform and centred components are orthogonal. Formula \eqref{nuc11:eq:5} is an exact inverse on compatible data; \eqref{nuc11:eq:8} specifies the stable extension used for noisy data, which need not satisfy the compatibility conditions of \eqref{nuc11:eq:5}.

\subsubsection{Example using the expressed register}

The register in \eqref{eq:k-vector} is

\begin{equation}
K=(234,543,140,729,659,824,621,58,914,794,146,601).
\label{nuc11:eq:10}
\end{equation}

It is an output of the preceding generator, used as an application case. The coefficients and readers are fixed before this register is evaluated. To store the memories using rational arithmetic, write \(z_i=m_i/\sqrt8\). The two effective registers are

\[
9z_0=(495,116,684,-108,-601,90,173,88,-313,-560,397,-461),
\]

\begin{equation}
81z_1=(2729,2503,3818,-5843,2583,-920,-4051,4338,-5312,-1583,233,1505).
\label{nuc11:eq:11}
\end{equation}

Substituting these registers into the reconstruction formulas, with \(q=6263\), reproduces the twelve components of \(K\). With \(q=0\), it reproduces \(K-(6263/12)\mathbf1\). The supplementary certificate first reconstructs the register from the memories and then compares it with \eqref{eq:k-vector}.

Adding a third memory \(D_3T_4T_3K\) raises the smallest positive eigenvalue of the Gram matrix to \(8/27\). The optimal amplification decreases from \(9/4\) to \(\sqrt{27/8}\). A third memory \(D_4T_4T_3K\) retains \(16/81\). The choice between repeating step \(3\) or step \(4\) has a stability effect calculable before evaluating the register.

\subsection{Reconstruction of the dimensional projector from two memories}

\subsubsection{Composition with the incidence already constructed}

The exceptional chart in \cref{exc:k-direccion} defines the orthogonal representation of \(A_5\) on the twelve classes of \(A_5/C_5\), with explicit representatives and character. Its three-dimensional isotypic projector is \eqref{eq:k-proyectores-incidencia}:

\begin{equation}
P_3=\frac3{60}\sum_{a\in A_5}\chi_3(a^{-1})\rho(a),
\quad P_3P_{11}=P_3,\quad P_3\mathbf1=0.
\label{nuc11:eq:12}
\end{equation}

The generated register precedes this representation. The result in \cref{lem:k-sector-no-nulo}, together with \eqref{eq:k-norma-direccion} and \eqref{eq:k-proyector-diez}, establishes

\[
u_K=P_3P_{11}K,\qquad
\|u_K\|^2=\frac{6638585+2275584\sqrt5}{20}>0,
\]

\begin{equation}
P_{10}(K)=P_{11}-\frac{u_Ku_K^*}{\|u_K\|^2}.
\label{nuc11:eq:13}
\end{equation}

The new composition of \eqref{nuc11:eq:8}, \eqref{nuc11:eq:12}, and \eqref{nuc11:eq:13} is

\begin{equation}
\boxed{
u_K=P_3L(m_0,m_1),\qquad
P_{10}(K)=P_{11}-
\frac{P_3L(m)\,[P_3L(m)]^*}{\|P_3L(m)\|^2}.
}
\label{nuc11:eq:14}
\end{equation}

\begin{proof}
The identities \(LM=P_{11}\) and \(P_3P_{11}=P_3\) give \(P_3LMK=P_3K=u_K\). Nonvanishing is evaluated in \eqref{eq:k-norma-direccion}. The rank-one quotient therefore coincides with that of \eqref{eq:k-proyector-diez}. The composition recovers the complete rank-ten orthogonal projector from the two memories, independently of the charge \(q\) and terminal reading \(T_4T_3K\).
\end{proof}

The reduction \(12\to11\) removes the uniform direction, and the reduction \(11\to10\) uses the vector selected by \(K\). The composition recovers the full projector matrix on the linear incidence carrier. Its dimensional realizations retain the maps of \cref{prop:k-reduccion-diez}.

\subsubsection{Projector error}

Let \(\varepsilon\) be the joint norm of the error in the two memories, and let \(\widehat u=P_3L(m+\delta m)\). Then \(\|\widehat u-u_K\|\le9\varepsilon/4\). If \(9\varepsilon/4<\|u_K\|\), the vector \(\widehat u\) is nonzero and

\begin{equation}
\boxed{
\|\widehat P_{10}-P_{10}\|
\le\frac{9\varepsilon}{4\|u_K\|}.
}
\label{nuc11:eq:15}
\end{equation}

\begin{proof}
The operator-norm distance between two orthogonal projectors onto lines is the sine of the angle between them. The distance from \(u_K\) to the line of \(\widehat u\) is at most \(\|u_K-\widehat u\|\). Division by \(\|u_K\|\) proves the bound. Here \(\|u_K\|\approx765{,}733\), and the coefficient is approximately \(0{,}00293836\).
\end{proof}
For registers close to \(\ker P_3\), stability of the direction depends on \(\|u_K\|\). The bound \(9/4\) for the centred register remains valid.

\subsubsection{A scalar descriptor of K, with explicit invariances}

The register \(K\), its scalar readings, and its invariants are distinguished. For \(K\notin\RR\mathbf1\), define

\begin{equation}
\eta_K=\frac{\|P_3P_{11}K\|^2}{\|P_{11}K\|^2}.
\label{nuc11:eq:16}
\end{equation}

For the register \eqref{eq:k-vector}, \(\sum_iK_i^2=4251257\) and \(\|P_{11}K\|^2=11789915/12\). Thus,

\begin{equation}
\eta_K=
\frac{3(6638585+2275584\sqrt5)}{5\cdot11789915}
=\frac{13089003+13653504\varphi_{\rm HMT}}{58949575}
\approx0,5967954228259091.
\label{nuc11:eq:17}
\end{equation}

The second expression uses the subsequent quadratic recognition of \(\varphi_{\rm HMT}\). We have \(0\le\eta_K\le1\). The ratio is invariant under \(K\mapsto aK+b\mathbf1\), with \(a\ne0\), and under simultaneous orthogonal changes of \(K,P_{11},P_3\). In the fixed chart, it is invariant under \(\rho(A_5)\), since both projectors commute with that representation. It is recoverable from the memories because both norms defining it are recoverable. It measures the quadratic fraction of the centred register lying in the chosen isotypic sector. The complete register also preserves the information identified by this scalar descriptor. The provenance of this additional formalization is documented in the supplement; the affine invariances proved here belong to the linear realization.

\subsection{Application to an effective vector of the exceptional construction}

The construction in \cref{km:orden-tres-reticular,km:estabilidad-vecino} uses twelve \(A_2\) fibres, their gluing, and the lattice neighbour. The tritic word is \(c^\star=(1,1,1,1,1,1,2,1,1,1,1,1)\). Its blocks act by \(C^{-1}\) when the component is \(1\) and by \(C\) when it is \(2\), in the charts specified by \cref{km:palabra-total,km:isometria-ternaria}. The resulting operator \(g\) satisfies \(g^3=I\), \(g^2+g+I=0\), and is orthogonal for the sum of the \(A_2\) metrics. Descent to the gluing and the neighbour is proved in \cref{km:estabilidad-vecino}. The following calculations act on its real carrier of dimension \(24\).

In each fibre,
\[
 C=\begin{pmatrix}0&-1\\1&-1\end{pmatrix},
 \qquad
 G_{A_2}=\begin{pmatrix}2&-1\\-1&2\end{pmatrix}.
\]
On the sum of the twelve fibres, define

\[
T=\frac{8I+g}{9},\qquad D=\frac{\sqrt8}{9}(g-I).
\]

With the adjoint relative to that metric, \(g^*=g^{-1}=-g-I\), and we obtain

\begin{equation}
T^*T=\frac{19}{27}I,
\qquad D^*D=\frac8{27}I,
\qquad D^{-1}=-\frac3{\sqrt8}(g+2I).
\label{nuc11:eq:18}
\end{equation}

The inverse comes from \((g-I)(g+2I)=-3I\). The neighbour construction fixes the vector in \eqref{exc:vector}, \(v=4\rho_1+\rho_2+\cdots+\rho_{12}\), with a root \(\rho_b\) of squared norm \(2\) in each fibre. Therefore,

\begin{equation}
\|v\|^2=54,\qquad \|Tv\|^2=38,\qquad \|Dv\|^2=16,
\qquad v=-\frac3{\sqrt8}(g+2I)Dv.
\label{nuc11:eq:19}
\end{equation}

A single complement recovers this complete vector. The reading preserves \(38\) units of squared norm and the register contains \(16\); their sum is \(54\). The result uses the vector in \eqref{exc:vector} and its lattice fibres.

\subsection{Conservation law in every exterior degree}

\subsubsection{Statement and proof}

Let \(P_0=I\), \(P_n=T_{n-1}\cdots T_0\), and

\[
\mathcal J_Nv=(P_Nv,D_0P_0v,\ldots,D_{N-1}P_{N-1}v).
\]

The law in \cref{nuclear:8,nuclear:10} preserves the form \(G_0\) through the orthogonal sum of the corresponding forms: \(\mathcal J_N^*G_{\rm sal}\mathcal J_N=G_0\). For each admissible finite exterior degree \(p\), it follows that

\begin{equation}
\boxed{
(\Lambda^p\mathcal J_N)^*(\Lambda^pG_{\rm sal})
(\Lambda^p\mathcal J_N)=\Lambda^pG_0.
}
\label{nuc11:eq:20}
\end{equation}

The induced form is defined by

\[
\langle x_1\wedge\cdots\wedge x_p,
y_1\wedge\cdots\wedge y_p\rangle_{\Lambda^pG}
=\det(\langle x_i,y_j\rangle_G)_{i,j=1}^p.
\]

\begin{proof}
Substitute the degree-one balance into each entry of this determinant. The equality holds on decomposable vectors and extends bilinearly because these span the exterior power.
\end{proof}

The relevant decomposition is

\[
\Lambda^p\Bigl(\bigoplus_aH_a\Bigr)
\simeq\bigoplus_{\sum n_a=p}\bigotimes_a\Lambda^{n_a}H_a.
\]

In the positive case, if \(\mathcal J_{p,\mathbf n}\) denotes the components of that map,

\begin{equation}
\boxed{\|\xi\|^2=\sum_{\sum n_a=p}
\|\mathcal J_{p,\mathbf n}\xi\|^2.}
\label{nuc11:eq:21}
\end{equation}

The purely visible, purely registered, and mixed components are preserved. For \(p=2\) and \(p=3\), the squared area and volume norms, respectively, are preserved. For an alternating or indefinite form, the exterior identity retains its bilinear interpretation.

\subsubsection{Exact distribution among the exterior sectors}

In the preceding Coxeter realization, \(a=19/27\) and \(b=8/27\). Normalizing \(T\) and \(D\) gives two isometries, so the component with \(k\) directions in memory has squared norm

\begin{equation}
\boxed{
\|\mathcal J_{p,k}\xi\|^2=
\binom pk a^{p-k}b^k\|\xi\|^2.
}
\label{nuc11:eq:22}
\end{equation}

This is proved by exterior expansion of \(x\mapsto(\sqrt a\,x,\sqrt b\,x)\). Components with distinct memory positions are orthogonal; the normalized isometries \(T\) and \(D\) preserve their norms.

For a unit-norm bivector,

\begin{center}
\begin{tabular}{@{}lr@{}}
\toprule
Sector & Squared norm\\
\midrule
Two directions in the reading & \(361/729\)\\
One direction in the reading and one in memory & \(304/729\)\\
Two directions in memory & \(64/729\)\\
\bottomrule
\end{tabular}
\end{center}

The sum is one. The mixed sector contains exactly \(304/729\), approximately \(41{,}7\%\), of the squared-area identity. In degree three, the four components are \((6859,8664,3648,512)/19683\); the mixed components sum to \(12312/19683\).

The law applies to \(0\le p\le24\) on the carrier of the twelve \(A_2\) fibres, and functorially to other generated carriers. Exterior degree, refinement resolution, the scale of a metric, and physical dimension have distinct typings. The change \(G\mapsto s^2G\), with \(s>0\), multiplies the squared norm of degree \(p\) by \(s^{2p}\); both sides of the exterior identity transform in the same way. The realization of dimensions \(10\) and \(11\) preserves the previously specified dimensional maps.

\subsubsection{Naturality and arbitrary depth}

The generated refinement of histories preserves measure through

\[
R\delta_h=\sum_{\epsilon\in\operatorname{Out}(h)}
\sqrt{p_h(\epsilon)}\,\delta_{h\epsilon}.
\]

The probabilities come from the multiplicities and normalization of the APP--TRIT--TPK emissions in \cref{iv:cont:medida,nuclear:1}. The register identifies the preceding state of each emission. If the transport square is \(\mathcal J'_NR=(\bigoplus_aR_a)\mathcal J_N\), then

\begin{equation}
(\Lambda^pJ'_N)(\Lambda^pR)
=\Lambda^p(\bigoplus R_a)(\Lambda^pJ_N).
\label{nuc11:eq:23}
\end{equation}

Apply \(\Lambda^p(AB)=\Lambda^p(A)\Lambda^p(B)\). The squares of the five correlative operations and their common limit are those of \cref{iv:rectangular-natural,iv:cont:memoria-natural,iv:cont:mapa-correlativo}; all act within the same construction of the continuum.

In the nonstationary positive case, \cref{nuclear:10} proves \(P_N^*P_N\downarrow A_\infty\) strongly, and

\[
\mathcal Jv=(A_\infty^{1/2}v,(D_nP_nv)_{n\ge0}),
\qquad \mathcal J^*\mathcal J=I.
\]

By density, \(\Lambda^p\mathcal J\) is isometric for each finite degree \(p\). The exterior sum extends to finitely supported multi-indices over the countable Hilbert direct sum; Parseval justifies passage to the limit. The terminal component \(A_\infty^{1/2}\) and the memories jointly preserve the initial form at arbitrary depth.

\subsubsection{Precision of area and volume reconstruction}

For the reader \(M\), order the eleven positive eigenvalues of \(G\), with multiplicity, as \(\lambda_1\le\cdots\le\lambda_{11}\). In degrees \(1\le p\le11\) of \(\mathbf1^\perp\), the singular-value decomposition induces

\begin{equation}
\|(\Lambda^pM)^\dagger\|
=(\lambda_1\cdots\lambda_p)^{-1/2}.
\label{nuc11:eq:24}
\end{equation}

The exterior modes use \(p\) distinct directions. The computed spectrum gives the optimal constants \(9/4\), \(81/16\), \(243\sqrt6/64\), and \(2187/128\) for \(p=1,2,3,4\). In degree twelve, \(\Lambda^{12}M=0\); recovery of the complete volume also incorporates the uniform charge. The tensor product of \(p\) independent readers on the zero-mean subspace has bound \((9/4)^p\). Tensoring the reader with an identity preserves the bound \(9/4\). The complete register \(\mathcal J\) has an inverse of norm one on its image in every degree.

\subsection{Transport memory and relative holonomy}

The reversible coupling in \cref{nuclear:10} is

\[
U(B,C)=\frac19
\begin{pmatrix}8B+C&\sqrt8(C-B)\\
\sqrt8(C-B)&B+8C\end{pmatrix}
=Q^*\operatorname{diag}(B,C)Q,
\]

with \(Q\) constant and orthogonal, fixed by the partition \(8+1\). Hence \(U(B_2,C_2)U(B_1,C_1)=U(B_2B_1,C_2C_1)\), preserving the order of factors. On a loop \(\gamma\), let \(H_B,H_C\) be the ordered products of the two transports. The memory block of the complete evolution is

\[
D_\gamma=\frac{\sqrt8}{9}(H_C-H_B).
\]

We reconstruct exactly

\begin{equation}
\boxed{H_B^{-1}H_C=I+\frac9{\sqrt8}H_B^{-1}D_\gamma.}
\label{nuc11:eq:25}
\end{equation}

The equality is an operator equality: recording the block's response on a basis recovers the complete relative holonomy. Its evaluation on a vector recovers the corresponding action. Under a chart change at the base point, both sides are conjugated. If \(H_B=I\), the complement determines \(H_C-I\). This relates memory and transport defect on a circuit; the differential realization additionally uses its connection and the limit of circuits.

\subsubsection{An elliptic–parabolic loop produces hyperbolic self-scaling}

The construction in \cref{euler-trit-generadores} contains the Coxeter operator \(C\) and the nilpotent
\[
 N_0=\begin{pmatrix}0&1\\0&0\end{pmatrix}
\]
on their respective carriers. Consider their realizations on a common oriented plane, using the indicated bases and the form
\[
 \Omega=\begin{pmatrix}0&1\\-1&0\end{pmatrix}.
\]
The Coxeter operator \(C\) and transport \(I+N_0\) have determinant one and preserve \(\Omega\). This chart allows them to be composed while preserving the distinction between the generators of their regimes and their own positive metrics.

Let
\[
 P=I+N_0=\begin{pmatrix}1&1\\0&1\end{pmatrix}.
\]
Since \(N_0^2=0\), we have \(P^n=I+nN_0\) for every \(n\in\ZZ\). The ordered loop of four transports gives

\begin{equation}
H_n=C^{-1}P^{-n}CP^n
=\begin{pmatrix}1+n&n^2\\n&n^2-n+1\end{pmatrix}.
\label{nuc11:eq:28}
\end{equation}

For \(n\ne0\), its normalized difference is

\[
F_n=\frac{H_n-I}{n}
=\begin{pmatrix}1&n\\1&n-1\end{pmatrix}.
\]

Multiplying these matrices gives, without introducing spectral values,

\begin{equation}
\boxed{F_n^2=H_n,\qquad F_n^2-nF_n-I=0.}
\label{nuc11:eq:29}
\end{equation}

With \(H_B=I\), memory is \(D_{\gamma,n}=(\sqrt8/9)nF_n\). Consequently,

\begin{equation}
\boxed{H_n=\left(\frac9{n\sqrt8}D_{\gamma,n}\right)^2.}
\label{nuc11:eq:30}
\end{equation}

The matrix-valued memory response reconstructs hyperbolic self-scaling. For \(n>0\), subsequent spectral recognition identifies

\begin{equation}
\rho_n=\frac{n+\sqrt{n^2+4}}2,\quad
\operatorname{Spec}F_n=\{\rho_n,-\rho_n^{-1}\},\quad
\operatorname{Spec}H_n=\{\rho_n^2,\rho_n^{-2}\}.
\label{nuc11:eq:31}
\end{equation}

The case \(n=1\) produces the golden polynomial. The relation to the operator \(F_{\rm av}=\begin{pmatrix}0&1\\1&1\end{pmatrix}\) generated by the calendar in \eqref{eq:fav} preserves the chart change:

\begin{equation}
F_1=PF_{\rm av}P^{-1},\qquad
H_1=PF_{\rm av}^{\,2}P^{-1}.
\label{nuc11:eq:32}
\end{equation}

The transport \(P\), of determinant one, intertwines the two matrices and preserves \(\Omega\). The case \(n=2\) produces the spectral readings \(1+\sqrt2\) and \(3+2\sqrt2\). The composition relates these realizations while preserving their respective genealogies; the selection of a four-factor word remains specified as part of the reader.

In the same example \(n=1\), the oriented balance is

\begin{equation}
D_\gamma^{\mathsf T}\Omega D_\gamma=-\frac8{81}\Omega,
\qquad
T_\gamma^{\mathsf T}\Omega T_\gamma=\frac{89}{81}\Omega.
\label{nuc11:eq:33}
\end{equation}

The sum of the two contributions is \(\Omega\). The memory term has opposite orientation: the visible contribution \(89/81\) is compensated by \(-8/81\). This alternating balance and the positive allocation \(19/27+8/27=1\) preserve their respective forms.

The expression for \(H_n\) holds for every \(n\in\ZZ\). The definition of \(F_n\) and reconstruction through its square use \(n\ne0\); the preceding spectral parametrization uses \(n>0\). In general,
\[
 D_{\gamma,n}^{\mathsf T}\Omega D_{\gamma,n}
 =-\frac{8n^2}{81}\Omega,\qquad
 T_{\gamma,n}^{\mathsf T}\Omega T_{\gamma,n}
 =\left(1+\frac{8n^2}{81}\right)\Omega.
\]
The identities \(N_0^2=0\) and \(C^2+C+I=0\) prove the complete collection; the supplementary program checks its exact specializations. The provenance of this additional formalization is documented in the supplement.

\subsection{Symmetry, action, and conserved charge}

Let \((E_n)_{n\in\ZZ}\) be finite-dimensional real vector spaces and \(U_n:E_n\to E_{n+1}\) the invertible transports of the complete evolution. The cotangent lift acts on \(q_n\in E_n\) and \(p_n\in E_n^*\). In dual bases, define the auxiliary discrete Lagrangian
\[
 L_n(q_n,q_{n+1},p_{n+1})
 =p_{n+1}^{\mathsf T}(q_{n+1}-U_nq_n).
\]
The action is used in a local sense: for each finitely supported variation, sum \(L_n\) over a finite interval containing all affected terms. The first variation is independent of any enlargement of that interval. Specifically,
\begin{equation}
 \delta\mathscr S
 =\sum_{n\in\ZZ}
 \left[
  \delta p_n^{\mathsf T}(q_n-U_{n-1}q_{n-1})
  +\delta q_n^{\mathsf T}(p_n-U_n^{\mathsf T}p_{n+1})
 \right],
 \label{nuc11:accion-local}
\end{equation}
where the sum is finite for every admitted variation. Stationarity with respect to all these variations is equivalent to
\[
 q_{n+1}=U_nq_n,\qquad
 p_n=U_n^{\mathsf T}p_{n+1}
 \quad(n\in\ZZ).
\]
The second equation is \(p_{n+1}=U_n^{-\mathsf T}p_n\), that is, the cotangent lift of transport.

\begin{proposition}[Variational charge of intertwined transport]
\label{nuc11:noether}
Let \(A_n\in\operatorname{End}(E_n)\) be a collection satisfying
\[
 A_{n+1}U_n=U_nA_n.
\]
On the stationary solutions of \eqref{nuc11:accion-local}, the charge
\begin{equation}
 \boxed{J_n=p_n^{\mathsf T}A_nq_n=J_{n+1}}
 \label{nuc11:carga}
\end{equation}
is constant.
\end{proposition}
\begin{proof}
The evolution equations and intertwining give
\[
 J_{n+1}
 =p_{n+1}^{\mathsf T}A_{n+1}U_nq_n
 =p_{n+1}^{\mathsf T}U_nA_nq_n
 =p_n^{\mathsf T}A_nq_n.
\]
Conservation also has an explicit variational derivation. For constant parameter \(\varepsilon\), the transformations
\[
 q_n\mapsto e^{\varepsilon A_n}q_n,\qquad
 p_n\mapsto e^{-\varepsilon A_n^{\mathsf T}}p_n
\]
preserve every \(L_n\), because \(e^{\varepsilon A_{n+1}}U_n=U_ne^{\varepsilon A_n}\). If the parameter is replaced by a finitely supported sequence \((\varepsilon_n)\), the first variation is
\[
 \delta L_n
 =(\varepsilon_{n+1}-\varepsilon_n)
   p_{n+1}^{\mathsf T}U_nA_nq_n.
\]
On a solution, the last factor is \(J_n\). Finite summation by parts gives
\[
 \delta\mathscr S
 =\sum_{n\in\ZZ}\varepsilon_n(J_{n-1}-J_n).
\]
Stationarity for every finitely supported sequence produces \eqref{nuc11:carga}.
\end{proof}

For the preceding cyclic readers, take
\[
 A_0=S-S^{-1},\qquad
 \mathcal A=\operatorname{diag}(A_0,A_0)
\]
on the complete carrier \(V\oplus V\). The matrix \(\mathcal A\) commutes with \(U(I,S^3)\) and \(U(I,S^4)\), since its blocks are polynomials in \(S\). Consequently, \(p_n^{\mathsf T}\mathcal A q_n\) preserves an oriented correlation between neighbouring positions.

The construction specifies a Lagrangian, its symmetry, and the corresponding conserved charge, in the variational sense developed in \cref{nuclear:7}. Its domain is the cotangent realization of the readers. The physical action and its units retain the operators and normalizations of the corresponding constructions. The auxiliary Lagrangian is obtained after the transport \(U_n\).
